\section*{Bab \ref{ch:g7:plate-tectonics} --- Tektonika Lempeng: Teka-teki yang Bergerak}

\begin{solution}{exo:g7:plate-tectonics:1}
\omterm{def:g7:plate-tectonics:continental-drift}{Pergeseran benua} adalah gagasan bahwa \omterm{def:g3:continents-and-oceans:continent}{benua-benua} bergerak perlahan di
permukaan Bumi, sehingga letaknya telah berubah dari zaman ke zaman.
Petunjuk-petunjuk Wegener: kecocokan garis pantai (Amerika Selatan dan
Afrika); fosil yang sama di \omterm{def:g3:continents-and-oceans:continent}{benua-benua} yang kini dipisahkan \omterm{def:g2:water-on-earth:ocean}{samudra};
\omterm{def:g3:stones-pebbles-sand:rock}{batuan} dan \omterm{def:g4:world-physical-map:range}{pegunungan} tua yang sama di kedua sisi \omterm{def:g2:water-on-earth:ocean}{samudra}; jejak \omterm{def:g5:weather-and-climate:climate}{iklim} purba
yang tidak sesuai dengan tempatnya (batu bara di Antarktika, endapan gletser
di Afrika yang panas dan kering).
\end{solution}

\begin{solution}{exo:g7:plate-tectonics:2}
Divergen: \omterm{def:g7:plate-tectonics:plate}{lempeng-lempeng} saling menjauh. Konvergen: \omterm{def:g7:plate-tectonics:plate}{lempeng-lempeng} saling
mendekat. Transform: \omterm{def:g7:plate-tectonics:plate}{lempeng-lempeng} saling bergeser berpapasan, berdampingan.
\end{solution}

\begin{solution}{exo:g7:plate-tectonics:3}
Misalnya, di sekeliling \omterm{def:g2:water-on-earth:ocean}{Samudra} Pasifik (Jepang, Indonesia, \omterm{def:g4:world-physical-map:range}{Pegunungan}
Andes), dan di sabuk yang membentang dari \omterm{def:g2:water-on-earth:ocean}{Laut} Tengah sampai Himalaya.
Hampir tidak ada \omterm{def:g4:volcanoes-earthquakes-first-look:earthquake}{gempa bumi} besar di tengah Afrika atau di tengah Australia
(juga di Amerika Utara sebelah timur \omterm{def:g4:world-physical-map:range}{Pegunungan} Rocky, atau di Brasil).
\end{solution}

\begin{solution}{exo:g7:plate-tectonics:4}
Hewan \omterm{def:g2:water-on-earth:salt-fresh}{air tawar} tidak mungkin menyeberangi \omterm{def:g2:water-on-earth:ocean}{samudra} \omterm{def:g2:water-on-earth:salt-fresh}{air asin} selebar ribuan
kilometer. Ditemukannya hewan itu di kedua sisi berarti bahwa kedua tempat
itu dulu dihubungkan oleh daratan, dalam satu \omterm{def:g3:continents-and-oceans:continent}{benua}, pada masa hidup
\emph{Mesosaurus}.
\end{solution}

\begin{solution}{exo:g7:plate-tectonics:5}
Batu bara terbentuk dari tumbuhan hutan rawa yang hangat dan basah: \omterm{def:g5:weather-and-climate:climate}{iklimnya}
hangat dan lembap. Hutan seperti itu tidak dapat tumbuh di dekat \omterm{def:g3:continents-and-oceans:poles}{Kutub
Selatan} pada masa kini, jadi Antarktika dulu pasti jauh lebih dekat ke
\omterm{def:g3:continents-and-oceans:equator}{khatulistiwa}: \omterm{def:g3:continents-and-oceans:continent}{benua} itu telah bergerak.
\end{solution}

\begin{solution}{exo:g7:plate-tectonics:6}
Punggung Tengah Atlantik: divergen. Andes: konvergen (\omterm{def:g7:plate-tectonics:plate}{lempeng} \omterm{def:g2:water-on-earth:ocean}{samudra}
menunjam ke bawah \omterm{def:g3:continents-and-oceans:continent}{benua}). Sesar San Andreas: transform. Himalaya: konvergen
(dua \omterm{def:g3:continents-and-oceans:continent}{benua} bertabrakan). Rekahan Afrika Timur: divergen (sebuah \omterm{def:g3:continents-and-oceans:continent}{benua} yang
sedang terbelah).
\end{solution}

\begin{solution}{exo:g7:plate-tectonics:7}
$2.5 \times 10 = 25$ kilometer per juta tahun; dalam 60 juta tahun:
$25 \times 60 = \qty{1500}{km}$.
\end{solution}

\begin{solution}{exo:g7:plate-tectonics:8}
$\qty{7}{m} = \qty{700}{cm}$, dan $700 \div 9 \approx 78$: kira-kira
\qty{78}{cm} per tahun, sekitar 30 kali rata-rata \qty{2.5}{cm} per tahun.
\omterm{def:g7:plate-tectonics:plate}{Lempeng-lempeng} tidak saling menjauh secara mulus di permukaan: tegangan
menumpuk selama bertahun-tahun lalu dilepaskan dalam letupan-letupan
singkat, ketika retakan-retakan terbuka bermeter-meter sekaligus. Dalam
jangka panjang, tahun-tahun yang tenang dan letupan-letupan itu dirata-ratakan
menjadi beberapa sentimeter per tahun.
\end{solution}

\begin{solution}{exo:g7:plate-tectonics:9}
Ia tidak dapat mengatakan gaya apa yang menggerakkan \omterm{def:g3:continents-and-oceans:continent}{benua}, dan gerak yang
dibayangkannya (\omterm{def:g3:continents-and-oceans:continent}{benua} membajak dasar \omterm{def:g2:water-on-earth:ocean}{samudra}) mustahil. Dalam \omterm{def:g7:plate-tectonics:plate-tectonics}{tektonika
lempeng}, \omterm{def:g3:continents-and-oceans:continent}{benua-benua} tidak membajak apa pun: \omterm{def:g3:continents-and-oceans:continent}{benua} itu dibawa, bersama dasar
\omterm{def:g2:water-on-earth:ocean}{samudra} di sekelilingnya, oleh \omterm{def:g7:plate-tectonics:plate}{lempeng-lempeng} yang bergerak sebagai
kepingan utuh cangkang luar Bumi.
\end{solution}

\begin{solution}{exo:g7:plate-tectonics:10}
Paling tinggi di punggung \omterm{def:g2:water-on-earth:ocean}{samudra}, paling dalam di palung. \omterm{def:g7:plate-tectonics:plate}{Lempeng} baru
dibuat di punggung \omterm{def:g2:water-on-earth:ocean}{samudra}, tempat \omterm{def:g3:stones-pebbles-sand:rock}{batuan} panas naik dari \omterm{def:g6:inside-the-earth:layers}{mantel}; \omterm{def:g7:plate-tectonics:plate}{lempeng}
\omterm{def:g2:water-on-earth:ocean}{samudra} kembali ke dalam \omterm{def:g6:inside-the-earth:layers}{mantel} di palung, tempat \omterm{def:g7:plate-tectonics:plate}{lempeng} itu menunjam ke
bawah \omterm{def:g3:continents-and-oceans:continent}{benua}.
\end{solution}

\begin{solution}{exo:g7:plate-tectonics:11}
$14.2 \div 2.2 \approx 6.45$: kira-kira 6 kali lebih cepat.
Dalam kilometer per juta tahun, kecepatannya 142 dan 22. Waktu untuk
\qty{1000}{km}: $1000 \div 142 \approx 7$ juta tahun untuk Punggung
Pasifik Timur; $1000 \div 22 \approx 45$ juta tahun untuk punggung Atlantik
Utara.
\end{solution}

\begin{solution}{exo:g7:plate-tectonics:12}
(a) $6472 - 4578 = 1894$ \omterm{def:g4:volcanoes-earthquakes-first-look:earthquake}{gempa bumi}. (b) $4578 \div 7031 \approx 0.651$,
yaitu 65\,\%. (c) Garis pada \omterm{def:g2:drawing-to-map:map}{peta} menandai palung, tempat \omterm{def:g7:plate-tectonics:plate}{lempeng} \omterm{def:g2:water-on-earth:ocean}{samudra}
mulai menunjam. \omterm{def:g7:plate-tectonics:plate}{Lempeng} itu tidak turun tegak lurus: \omterm{def:g7:plate-tectonics:plate}{lempeng} itu menunjam
miring ke bawah \omterm{def:g7:plate-tectonics:plate}{lempeng} yang lain, makin lama makin dalam saat masuk ke
bawahnya. \omterm{def:g4:volcanoes-earthquakes-first-look:earthquake}{Gempa bumi} terjadi di sepanjang \omterm{def:g7:plate-tectonics:plate}{lempeng} yang menunjam, sehingga
banyak di antaranya terletak di bawah \omterm{def:g7:plate-tectonics:plate}{lempeng} yang menindih, pada sisi palung
yang sama, puluhan atau ratusan kilometer dari garis itu. \omterm{def:g7:plate-tectonics:plate-boundary}{Batas lempeng}
adalah sebuah zona, bukan garis tipis.
\end{solution}

\begin{solution}{pb:g7:plate-tectonics:1}
\textbf{1.} \emph{Mesosaurus}, \emph{Cynognathus}, dan \emph{Glossopteris}
(pita \emph{Lystrosaurus} menghubungkan Afrika, India, dan Antarktika).

\textbf{2.} Hewan itu hidup di \omterm{def:g2:water-on-earth:salt-fresh}{air tawar}: hewan itu tidak mungkin
menyeberangi \omterm{def:g2:water-on-earth:ocean}{samudra} \omterm{def:g2:water-on-earth:salt-fresh}{air asin} selebar ribuan kilometer.

\textbf{3.} Tonjolan Brasil bagian timur laut cocok dengan lekukan Teluk
Guinea, di pantai barat Afrika.

\textbf{4.} \omterm{def:g7:plate-tectonics:plate-boundary}{Batas divergen}, yaitu Punggung Tengah Atlantik: \omterm{def:g7:plate-tectonics:plate}{lempeng} Amerika
Selatan bergerak ke barat dan \omterm{def:g7:plate-tectonics:plate}{lempeng} Afrika ke timur, saling menjauh.

\textbf{5.} Kedua pantai itu terletak pada sisi yang berlawanan dari
\omterm{def:g6:latitude-longitude:longitude}{meridian utama}: $43.2 + 14.4 = \ang{57.6}$.

\textbf{6.} $57.6 \times 102 = 5875.2$: kira-kira \qty{5880}{km}.

\textbf{7.} \omterm{def:g6:latitude-longitude:longitude}{Meridian-meridian} makin rapat jika kita bergerak dari
\omterm{def:g3:continents-and-oceans:equator}{khatulistiwa} ke arah \omterm{def:g3:continents-and-oceans:poles}{kutub}. \omterm{def:g6:latitude-longitude:latitude}{Paralel} \ang{23}\,LS adalah lingkaran yang
lebih kecil daripada \omterm{def:g3:continents-and-oceans:equator}{khatulistiwa}, sehingga setiap derajat dari 360
derajatnya lebih pendek.

\textbf{8.} Kira-kira separuh lebarnya: $5880 \div 2 = \qty{2940}{km}$.

\textbf{9.} $3.4 \times 10 = 34$ kilometer per juta tahun.

\textbf{10.} $5880 \div 34 \approx 173$: kira-kira 170 juta tahun.

\textbf{11.} Ya: 170 juta tahun terletak antara 100 dan 200 juta tahun, dan
waktu itu datang sesudah awal terpecahnya \omterm{def:g7:plate-tectonics:supercontinent}{Pangea}, 225 sampai 200 juta tahun
yang lalu. Perhitungan sederhana pada \omterm{def:g2:drawing-to-map:map}{peta} memberikan orde besaran yang
benar.

\textbf{12.} Misalnya: kecepatannya tidak selalu sama sepanjang waktu;
\omterm{def:g3:continents-and-oceans:continent}{benua-benua} itu meregang dan menipis sebelum terbelah, sehingga pantai masa
kini tidak persis sama dengan garis tempat \omterm{def:g3:continents-and-oceans:continent}{benua} itu patah.
\end{solution}
